%% Complete derivations for the Ciardelli 2023 propositional selection.
\begin{annexitem}{annex:p8.inq.bases}{thm:p08-inq-bases}{Minimal generators and normal propositions}
\begin{proof}
Suppose a generator \(T\) is independent, and remove a member \(a\).
No remaining member contains \(a\), so \(a\) no longer belongs to the
downward closure. Thus every proper subfamily fails to generate the same
proposition. Conversely, if \(a\subsetneq b\) are in \(T\), removing \(a\)
does not change its downward closure: every subset of \(a\) is already a
subset of \(b\). This proves minimality exactly when independent.

If \(\mathcal A\) is normal, its alternatives generate it and are independent,
so form a basis. Let \(T\) be any basis. Every alternative \(a\) is contained
in a member \(b\in T\subseteq\mathcal A\); maximality forces \(a=b\).
Hence all alternatives belong to \(T\). To prove that no other members
occur, suppose \(a\in T\) is not maximal. Choose \(b\in\mathcal A\) properly
containing it, and then \(c\in T\) containing \(b\). This gives
\(a\subsetneq c\), contradicting independence. Thus
\(T=\operatorname{Alt}(\mathcal A)\), and normality follows from its being
a generator. This proves existence and uniqueness of a basis together.
Finally, finite subsets of an infinite set form the stated counterexample:
any finite subset can be extended by one more element, so none is maximal.
\end{proof}
\end{annexitem}

\begin{annexitem}{annex:p8.inq.expressive}{thm:p08-inq-expressive}{Finite profiles, expression and decision}
\begin{proof}
Support of a formula depends only on the profiles of its occurring atoms.
Prove this by induction, simultaneously across models. Atoms, falsum and
the two conjunction/disjunction clauses are immediate. For implication,
every subprofile of a state's profile is realized by a substate: retain
exactly the worlds with profiles in that subprofile. Thus the substate
tests match in any two states with the same profile. This proves the
induction and finite determination. Persistence and empty-state support
give the two inquisitive conditions.

Conversely let a property be \(Q\)-determined and inquisitive, with \(Q\)
finite. Evaluate it in the standard model \(\Omega_Q\), extending atoms
outside \(Q\) by false when needed. It determines a nonempty downward-closed
family \(\mathcal A\) of profile sets. A state supports \(\chi_A\) exactly
when its profile is included in \(A\): each \(\chi_v\) singles out one
profile, and \(\chi_A\) is classical. The finite inquisitive disjunction
over \(A\in\mathcal A\) therefore supports precisely the states whose
profile lies in \(\mathcal A\). This is the original property by finite
determination and downward closure. The family is nonempty because it
contains the empty profile; \(Q=\varnothing\) and empty states are included
by the empty-conjunction and empty-disjunction conventions.

Over finite \(Q\), there are only finitely many such families and hence
finitely many equivalence classes. If \(|Q|=n\), the standard space has
\(2^n\) worlds. Maximal members give a bijection between its nonempty
downward-closed families and the nonempty antichains in its subset lattice,
so the exact count is \(D(2^n)-1\). Flat families are precisely the
principal families generated by one subset of that space, giving
\(2^{2^n}\) flat classes. Over one atom, the standard space has
two worlds. Its nonempty downward-closed families are exactly those
generated by \(\varnothing\), either one singleton, the two singletons,
or the whole two-world set. These give the five listed formulas.
For two atoms the standard space has four worlds. Count antichains of its
subset lattice, since every nonempty downward-closed family has a unique
antichain of maximal members. The antichains using either the empty subset
or the full subset are the two corresponding singleton antichains. In the
three middle ranks, count first the chosen one-element subsets.
With none chosen, selecting zero, one, two, three or four three-element
subsets leaves respectively six, three, one, zero or zero available
two-element subsets. The count is
\(64+4\cdot8+6\cdot2+4+1=113\).
With one chosen singleton there are four choices: without a chosen triple
there are three available pairs, while with a triple only its complementary
triple is permitted and no pair remains, giving \(4(8+1)=36\).
With two chosen singletons there are six choices and one optional pair,
giving 12. With three there are four choices, and with four one choice.
Thus there are \(113+36+12+4+1+2=168\) antichains, including the empty
antichain. Excluding that one leaves 167 inquisitive propositions.

Logical split with no premises says that a formula is valid exactly when
one of its finitely many classical resolutions is valid. Classical truth
tables decide these tests. Alternatively, every model's relevant state
has a profile state in the finite standard model. If the whole standard
space supports a formula, persistence supports every such profile state,
and the profile induction supports every corresponding state in every
model. The converse is immediate. Finally, the implication/deduction
property turns finite-premise entailment into the stated validity problem.
All procedures require finite input; no claim about effectively processing
an arbitrary infinite premise set is made.
\end{proof}
\end{annexitem}

\begin{annexitem}{annex:p8.inq.substitution}{thm:p08-inq-substitution}{Substituting statements and substituting questions}
\begin{proof}
The valid entailment \(\neg\neg p\models p\) becomes invalid when
\(?p\) is substituted for \(p\), as witnessed by the two-world state
in the body. Thus unrestricted substitution fails.

For a substitution \(s\) mapping every atom to a flat formula, fix
\(M=(W,V)\) and define \(V_s(w,p)=1\) exactly when
\(M,w\models s(p)\). In \(M_s=(W,V_s)\), support of an atom at a
state \(t\) is equivalent to support of \(s(p)\) in \(M\), precisely
because \(s(p)\) is flat. Induction through every connective gives
\(M_s,t\models\varphi\) exactly when \(M,t\models s(\varphi)\).
For implication the induction is applied at every substate; the other
clauses follow directly. A counterstate to the substituted consequence
would therefore be a counterstate to the original one in \(M_s\).
This proves preservation of consequence under legitimate substitution.

If atoms can denote arbitrary downward-closed families, the same semantic
substitution construction instead works for arbitrary formulas, but an
atomic question then need not satisfy double-negation elimination. Its
atomic resolution can no longer be the singleton classical formula
\(\{p\}\). These concrete differences show that this is another system,
not the same \(\mathsf{InqB}\) with an additional admissible substitution rule.
\end{proof}
\end{annexitem}

\begin{annexitem}{annex:p8.inq.connectives}{thm:p08-inq-connectives}{Elimination and all connective-independence cases}
\begin{proof}
A uniform defining context gives elimination by induction on formulas,
replacing its connective at each occurrence after eliminating it in the
operands. In a substitution-closed logic, conversely, an elimination of
the formula with distinct atoms in its argument places supplies a defining
context by uniform substitution. We will not use that converse for InqB.

\paragraph{Elimination in the original signature.}
A falsum-free formula is true at a singleton where every atom is true,
by induction, so cannot express \(\bot\). A formula without inquisitive
disjunction is classical and flat, so cannot express \(?p\). At singletons,
formulas using only \(\bot,\land,\mathbin{\vee\!\vee}\) are monotone
Boolean functions of the atoms, so cannot express \(\neg p\); thus
implication is not eliminable either.
To eliminate conjunction, first replace each classical conjunction
\(\alpha\land\beta\) recursively by \(\neg(\alpha\to\neg\beta)\).
This preserves classical truth and hence support. Now put an arbitrary
formula into resolution normal form and perform that replacement in
each classical resolution. The resulting inquisitive disjunction is
conjunction-free and equivalent to the original formula.

\paragraph{Conjunction has no defining context.}
Let \(C(p_1,p_2)\) be any proposed conjunction-free context, allowing
other atoms. Choose fresh atoms \(p,q\) and substitute \(?p,?q\) in
the argument places. Use the four-world standard model for \(p,q\),
with all other atoms false. Every substituted subcontext is either
supported only at the empty state, or supported at every singleton and
at some two-world state. The induction is as follows. The two argument
places give \(?p\) and \(?q\), each having the latter property; falsum
and the remaining atoms give the former. For implication, if both
operands are empty-state-only the result is \(\top\). If only the
consequent is empty-state-only, negation of an antecedent true at all
singletons is empty-state-only. If the consequent has the latter
property, the implication is supported at all its supporting states,
and in particular at all singletons and a two-world state. For
inquisitive disjunction, both operands empty-state-only give that case;
otherwise the supporting states of an operand with the latter property
are retained. This proves the invariant for every context. But
\(?p\land?q\) is supported at all singletons and at no two-world
state, so no such context defines conjunction. Together with the
non-eliminability results, this proves independence of the four primitives.

\paragraph{Elimination with primitive negation.}
Negation cannot be eliminated from the signature without falsum:
a singleton making every atom true satisfies all formulas made only
with \(\land,\to,\mathbin{\vee\!\vee}\), but falsifies \(\neg p\).
Without inquisitive disjunction, every remaining formula is flat,
so that connective is again indispensable.
For a classical formula recursively replace falsum by
\(\neg(p_0\mathbin{\vee\!\vee}\neg p_0)\), implication by
\(\neg\neg(\neg\alpha\mathbin{\vee\!\vee}\beta)\), and conjunction
by \(\neg(\neg\alpha\mathbin{\vee\!\vee}\neg\beta)\), using
already translated operands and a fixed atom \(p_0\). Every result is
flat, and its singleton truth conditions are the required classical
ones. Apply this to each resolution in normal form. The result uses
only negation and inquisitive disjunction. This proves joint elimination
of conjunction and implication in the stated variant.

The preceding conjunction-undefinability invariant also works when
negation is primitive: negating an empty-state-only formula gives
\(\top\), and negating a formula supported at every singleton gives
the empty-state-only case.

\paragraph{Implication has no defining context without implication.}
Again substitute \(?p,?q\) for the two argument places in a four-world
model, with any other atoms false. Put \(A=?p\), \(B=?q\),
\(C=A\land B\), and \(D=A\mathbin{\vee\!\vee}B\).
Every implication-free context using
\(\neg,\land,\mathbin{\vee\!\vee}\) is equivalent in this model to
one of \(\bot,C,A,B,D,\top\). To prove the induction, their support
families form the six-element lattice
\(\bot<C<A,B<D<\top\), with \(A,B\) incomparable.
Intersection and union of these families give respectively their
lattice meet and join, staying among the six in every case.
Negation sends \(\bot\) to \(\top\) and every other member to
\(\bot\), because every other member supports every singleton.
The base contexts lie in the set, so the induction is complete.

The implication \(?p\to?q\) is none of these six. It is not
\(\bot\), since it supports singletons; it is not \(\top\), since
two worlds agreeing on \(p\) and disagreeing on \(q\) refute it.
It supports the two diagonal pairs, refuting equality with \(C,A,B,D\)
as follows: a diagonal pair supports none of \(A,B,D\), and no pair
supports \(C\). Therefore no such context defines implication.
This completes independence in the primitive-negation variant. The
unique minimal complete subsets follow because each of their members
was shown non-eliminable and the listed subsets were shown sufficient.
\end{proof}
\end{annexitem}

\begin{annexitem}{annex:p8.inq.coherence}{thm:p08-inq-coherence}{Coherence bounds, the polar fragment and partitions}
\begin{proof}
We use the contrapositive characterization: every failure of support has
a failing substate of size at most \(n\). Persistence gives the other
direction. This also shows that \(n\)-coherence implies \(m\)-coherence
for every \(m\geq n\).

An atomic or falsum failure has a singleton witness. A conjunction
failure is a failure of a conjunct, giving the maximum bound. If
\(s\) fails \(\varphi\to\psi\), take \(t\subseteq s\) supporting
\(\varphi\) and failing \(\psi\); a small failing substate for
\(\psi\) still supports \(\varphi\) and thus fails the implication.
If \(s\) fails an inquisitive disjunction, choose small failing substates
for both operands. Their union has at most the sum of the sizes and
fails both: support of either there would persist to its failing
substate. Structural induction gives all displayed bounds and existence
of a finite degree.

For exact degree \(n\), take the inquisitive disjunction of \(n\)
distinct atoms. In a model with worlds \(w_1,\ldots,w_n\), make
\(p_i\) false exactly at \(w_i\). Every state of size at most \(n-1\)
omits some \(w_i\) and supports \(p_i\), while the whole state
supports no disjunct. Together with the upper bound this gives degree
\(n\); the case \(n=1\) is the atomic case.

Every formula of \(\mathsf{InqC}\) is 2-coherent. The preceding bounds
handle atoms, falsum, conjunction and implication. For \(?\varphi\),
suppose \(s\) supports neither \(\varphi\) nor its negation.
Some \(w\in s\) makes \(\varphi\) true. If some \(v\in s\) makes
it false, the pair \(\{w,v\}\) fails both \(\varphi\) and its
negation and hence fails the question. Otherwise every singleton in
\(s\) supports \(\varphi\), but a substate of size at most two fails
\(\varphi\) by induction. That substate is nonempty and cannot support
\(\neg\varphi\), because its worlds make \(\varphi\) true.
It too fails the question. This proves the required bound.

Conversely, let an inquisitive state property be \(Q\)-determined and
2-coherent. In the finite standard model, exclude each forbidden singleton
profile \(v\) by the classical formula \(\neg\chi_v\). Exclude each
forbidden two-element profile set \(\{u,v\}\) by
\[
 (\chi_u\lor\chi_v)\to?\chi_u.
\]
This formula is supported exactly when the state does not contain both
profiles: restrict to those two profiles and ask whether the first holds.
Conjoin the finitely many exclusion formulas; use \(\top\) if there
are none. The resulting \(\mathsf{InqC}\) formula accepts exactly those
states whose one- and two-element subprofiles are permitted. Coherence
and profile determination make this exactly the required property.
The degree-three formula already constructed proves strict expressiveness.

For a state property characterized by its nonempty substates of size
at most \(n\), the empty state satisfies the characterization vacuously.
If \(s\) has the property and \(t\subseteq s\), every tested small
substate of \(t\) is one of those of \(s\), so \(t\) also has it.
This proves the assertion about arbitrary coherent state properties.

For partitions, a conjunction of finitely many polar questions groups
worlds by their common vector of answers, so its maximal nonempty
supporting states form disjoint blocks covering the universe.
Conversely, a partitional formula depends on a finite atom set \(Q\).
Its alternatives in the standard \(Q\)-model are finitely many disjoint
profile blocks covering \(\Omega_Q\). Let \(\alpha_i\) be the classical
formula \(\chi_{A_i}\) for each block. Resolving all \(?\alpha_i\)
means knowing exactly which block contains the actual profile: in a
nonempty state at least one block must be true throughout, since the
blocks cover the space, and the answers for different blocks cannot
both be positive. Conversely a state contained in a block resolves
all these polar questions. Thus their conjunction has the same support
family in the standard model. Profile invariance extends the equivalence
to every model, including the empty one.
The blocks define an equivalence relation, and a state lies within a
single block exactly when every pair in it is equivalent. Conversely,
this pair condition for an equivalence relation gives precisely its
blocks as maximal nonempty supports. This proves the equivalent
relation-based characterization as well.
\end{proof}
\end{annexitem}

\begin{annexitem}{annex:p8.inq.tensor}{thm:p08-inq-tensor}{Tensor expression and failure of uniform definition}
\begin{proof}
If \(s=s_1\cup s_2\) witnesses tensor support and \(t\subseteq s\),
the split \(t=(t\cap s_1)\cup(t\cap s_2)\) witnesses support at
\(t\), by induction and persistence. The empty state splits into two
empty states. Profile determination is preserved as well: a split of
profile sets can be lifted by taking all worlds whose profiles occur
in either chosen part, and taking profiles of a world split gives a
profile split. The induction through tensor therefore preserves all
three properties. The expressive-completeness theorem supplies an
equivalent \(\mathsf{InqB}\) formula for each formula in the extension.

Suppose a context in the original four connectives defined tensor.
Choose fresh \(p,q\), use their full four-world model with all other
atoms false, and fill both context arguments by
\(K=?p\land?q\). This formula supports exactly the empty state and
singletons. By induction, every filled subcontext is equivalent to
one of \(\bot,K,\top\). These families form a chain and are closed
under conjunction and inquisitive disjunction. Their implication table
is also closed: a smaller or equal antecedent-to-consequent comparison
gives \(\top\); \(\top\to K\) gives \(K\); and either \(K\)
or \(\top\) implying \(\bot\) gives \(\bot\), since every singleton
supports the antecedent. These cases exhaust the table.
Yet \(K\otimes K\) supports precisely states of size at most two:
such states split into at most two singletons, and a union of two
singleton-or-empty supports has at most two worlds. This is none of
the three families, as a two-world and a three-world state demonstrate.
The proposed context therefore cannot define tensor.
\end{proof}
\end{annexitem}

\begin{annexitem}{annex:p8.inq.polar-tasks}{thm:p08-inq-polar-tasks}{Polar contents and their refinement order}
\begin{proof}
Write \(A=|\alpha|_M\). A state supports \(?\alpha\) exactly when
it lies in \(A\) or in \(W\setminus A\). The whole \(W\) does so
exactly when one of these sets is all of \(W\); this proves the
tautology condition. If both parts are nonempty, they are exactly the
two maximal supports. Equal polar contents therefore have the same
unordered pair of complementary blocks. This gives equality of truth
sets or complementation; the trivial cases \(A=\varnothing,W\)
satisfy the same description.

If one polar content entails another nontrivial one, each of its
nonempty blocks must lie within one of the latter's two blocks.
A trivial antecedent cannot entail a nontrivial question. With both
antecedent blocks nonempty, coverage of \(W\) forces one to occupy
each conclusion block, and then forces equality of both blocks. Thus
the questions are equivalent. Finally, on nine worlds every subset
determines a polar content and only complementary subsets determine
the same one. The complementation involution has no fixed subset
when \(W\ne\varnothing\), so the number is \(2^9/2=256\).
\end{proof}
\end{annexitem}

\begin{annexitem}{annex:p8.inq.exercise-inferences}{thm:p08-inq-exercise-inferences}{Explicit mixed inference derivations}
\begin{proof}
For \(?p\land?q\), project the two questions and eliminate their
inquisitive disjunctions in turn. In the \(p,q\) branch, conjunction
introduction gives \(p\land q\), hence its polar question. In each
other branch at least one of \(\neg p,\neg q\) is available; assume
\(p\land q\), project the contradicted conjunct, derive falsum and
discharge to obtain \(\neg(p\land q)\), hence the same polar
conclusion. The two disjunction eliminations discharge the four
answer branches.

For the second entailment, assume \(p\). The first premise yields
\(?q\to?r\). From the second premise \(\neg q\), introduce
\(?q\), then eliminate the implication to obtain \(?r\).
Discharge \(p\) to get \(p\to?r\). Since \(p\) is classical,
split yields exactly \((p\to r)\mathbin{\vee\!\vee}(p\to\neg r)\).

For the third entailment, assume \(?m\) and split into its \(m\)
and \(\neg m\) branches. In the \(m\) branch, use the classical
disjunction \(a\lor b\) to derive the classical conclusion \(b\):
the \(a\) branch contradicts \(a\to\neg m\), while the \(b\)
branch gives \(b\) directly. Introduce \(a\mathbin{\vee\!\vee}b\).
In the \(\neg m\) branch, the symmetric classical case analysis,
using \(b\to m\), yields \(a\) and the same inquisitive
disjunction. Eliminate \(?m\) and discharge it for the implication.
Every use of classical disjunction elimination has a classical
conclusion; the question conclusion is introduced afterward.

For the final entailment, assume \(?h\) and eliminate its two
answer branches. Under \(h\), assume \(m\); the first background
premise gives \(h\to b\), hence \(b\), hence \(?b\).
Discharge \(m\), giving \(m\to?b\). Under \(\neg h\),
assume \(m\), and temporarily assume \(b\). The second background
premise gives \(\neg h\to\neg m\), contradicting \(m\).
Discharge \(b\) to get \(\neg b\), hence \(?b\), and
discharge \(m\). Both branches now conclude \(m\to?b\),
so inquisitive disjunction elimination supplies it from \(?h\).
\end{proof}
\end{annexitem}

\begin{annexitem}{annex:p8.inq.generator-entailment}{thm:p08-inq-generator-entailment}{Entailment between generators}
\begin{proof}
Suppose the entailment holds. For any choice \(a_i\in T_i\), the state
\(a=\bigcap_i a_i\) supports each premise by downward closure. It therefore
supports \(\psi\), so lies below some \(b\in U\). Conversely, if every
such intersection lies below a member of \(U\), take a state supporting
all premises and choose \(a_i\in T_i\) containing it. That state is contained
in their intersection and hence in the resulting \(b\); it supports the
conclusion. For no premises, the only intersection is \(W\), and the same
argument is exactly validity relative to the model. For statements the
truth-support bridge gives the specified singleton generators. Substituting
these into the criterion gives both mixed cases in the theorem.
\end{proof}
\end{annexitem}

\begin{annexitem}{annex:p8.inq.basic}{thm:p08-inq-basic}{Basic properties of support}
\begin{proof}
Prove persistence and empty-state support together by structural induction.
For atoms, universal truth on a state persists to a subset and is vacuous
on the empty set. For falsum the claims are immediate. Conjunction uses
both induction hypotheses; inquisitive disjunction uses the supported
disjunct, which remains supported in the smaller state. For implication,
every substate of a smaller state is also a substate of the original one,
so its universal condition persists. The sole substate of the empty state
is empty, where the induction hypotheses support both operands; thus the
implication is also supported there.

Restriction of a model leaves atomic values on the restricted state unchanged.
Conjunction and inquisitive disjunction preserve this observation. The
substates quantified over by implication are literally the same sets before
and after restriction. Falsum again depends only on emptiness. This proves
locality by induction, including every substate of the restriction.

At \(\{w\}\), conjunction and inquisitive disjunction give their ordinary
truth conditions. For implication the only substates are \(\varnothing\)
and \(\{w\}\). The former always satisfies the implication condition;
the latter requires that the antecedent be false or the consequent true.
Falsum is false at a world and atoms have their assigned values. The
defined classical operations therefore have the usual truth conditions too.

If \(w\) belongs to a supporting state for \(\varphi\), persistence gives
support at \(\{w\}\), and hence truth at \(w\). Conversely, truth at
\(w\) supplies that singleton supporting state. This proves the union
formula for the truth set. Finally, \(s\models\neg\varphi\) means that
every substate supporting \(\varphi\) is empty. This is equivalent to
the absence of a singleton supporting \(\varphi\): a nonempty supporting
state would supply one by persistence. It is therefore equivalent to
\(s\cap|\varphi|_M=\varnothing\).
\end{proof}
\end{annexitem}

\begin{annexitem}{annex:p8.inq.flatness}{thm:p08-inq-flatness}{Classical support, negation and flatness}
\begin{proof}
For a classical formula prove \(s\models\alpha\) exactly when
\(s\subseteq|\alpha|_M\) by induction. Atoms and falsum are immediate,
as is conjunction. For \(\alpha\to\beta\), the induction hypotheses
turn the support condition into
\(s\cap|\alpha|_M\subseteq|\beta|_M\), which is exactly containment
in its material truth set. This proves flatness of the whole classical
fragment.

Replacing each inquisitive disjunction by classical disjunction leaves
singleton truth unchanged, by the basic truth clauses. Thus a formula
and its classical variant have the same truth set. A flat formula is
determined by this set and is consequently equivalent to its classical
variant. Conversely, equivalence to a classical formula transfers flatness.

Conjunction preserves flatness directly. If a state fails to support
\(\varphi\to\alpha\) with flat \(\alpha\), some substate \(t\)
supports \(\varphi\) but not \(\alpha\). Flatness supplies a world
\(w\in t\) where \(\alpha\) is false, while persistence makes
\(\varphi\) true there. This world falsifies the implication. Together
with persistence this proves its flatness. Apply the result to \(\bot\)
to get flatness of every negation, and hence of the defined classical
disjunction. Double negation is flat and shares the original formula's
truth set, so it is equivalent to the classical variant. A supporting
state has all its worlds in that truth set, so it supports the double
negation. The equivalence with double negation is therefore exactly flatness.

If a formula is flat, unions of supporting states still lie in its truth
set and support it, including the empty union. Conversely, if supporting
states are closed under arbitrary unions, a state all of whose singletons
support the formula is the union of those supporting singletons and so
supports it. This is flatness.
\end{proof}
\end{annexitem}

\begin{annexitem}{annex:p8.inq.implication}{thm:p08-inq-implication}{Context, supposition and statements}
\begin{proof}
The first assertion is the implication clause itself. For a nested
implication, induction on the finite number of antecedents uses two facts:
a substate supporting earlier premises retains them in any further substate,
and the current substate can itself be chosen in the next implication test.
These show that all the successive subset quantifiers amount precisely to
testing substates supporting every listed premise.

If \(\alpha\) is flat, substates of \(s\) supporting it are exactly the
subsets of \(s\cap|\alpha|_M\). Persistence makes support at all these
subsets equivalent to support at their largest member. For a generator
\(T\) of \([\varphi]_M\), each \(s\cap a\), \(a\in T\), supports
\(\varphi\). Hence support of \(\varphi\to\psi\) implies support of
\(\psi\) there. Conversely, any \(t\subseteq s\) supporting \(\varphi\)
lies below some \(a\in T\), and thus below \(s\cap a\); support of
\(\psi\) follows by persistence.

With a flat consequent, failure of implication is witnessed by one world,
as proved in the flatness lemma. Its truth condition depends only on the
antecedent's truth set, shared with its classical variant. Thus replacing
the antecedent by that variant preserves the implication. For consequence
with flat conclusion, necessity of the worldwise condition follows by
testing singleton states. For sufficiency, any state supporting all premises
has every premise true at every world by persistence. The worldwise
condition makes the conclusion true everywhere, and its flatness makes
it supported by the state. Classical variants preserve those premise truth
conditions. The assertion about presuppositions is the singleton-premise
instance of this argument.
\end{proof}
\end{annexitem}

\begin{annexitem}{annex:p8.inq.resolution-semantics}{thm:p08-inq-resolution-semantics}{All resolution clauses and semantic normal form}
\begin{proof}
Structural induction proves finiteness, nonemptiness and classicality.
For conjunction take the finite nonempty product of the resolution sets;
for inquisitive disjunction take their union. For implication, the finite
set of functions between two finite nonempty sets is nonempty, and each
corresponding conjunction is classical. The same induction for a classical
formula yields exactly that formula as the sole resolution, with the
specified singleton-conjunction convention.

We prove the support equivalence by induction. Atoms and falsum have
singleton resolution sets. A supporting conjunction supports each operand,
so by induction it supports some resolution of each and their conjunction.
Conversely a supporting resolution of the conjunction supports both selected
operand resolutions and hence both operands. For inquisitive disjunction,
support of one operand is equivalent to support of one resolution in the
union of their resolution sets, in either direction.

Suppose \(s\models\varphi\to\psi\). For every \(\alpha\in R(\varphi)\),
the state \(s\cap|\alpha|_M\) supports \(\alpha\), since it is classical,
and therefore supports \(\varphi\) by induction. It supports \(\psi\)
by the implication assumption, and hence some \(\beta\in R(\psi)\)
by induction. The flat-antecedent implication clause gives
\(s\models\alpha\to\beta\). Choose one such \(\beta=f(\alpha)\)
for each member of the finite set \(R(\varphi)\). The resulting
\(\gamma_f\) is a resolution of the implication supported at \(s\).
Conversely, suppose \(s\models\gamma_f\) and \(t\subseteq s\) supports
\(\varphi\). By induction it supports some \(\alpha\in R(\varphi)\).
The corresponding conjunct of \(\gamma_f\), restricted to \(t\), gives
support of \(f(\alpha)\), and induction gives support of \(\psi\).
This proves the implication case completely.

The finite inquisitive disjunction of resolutions has exactly the same
support condition. Flatness of resolutions makes their truth sets a
generator. In a finite generator, every member is included in a maximal
one: repeatedly taking a proper larger member must stop. Thus every
supporting state lies below an alternative, establishing normality of all
propositional formulas. A unique alternative generates a principal downward
set whose union is that alternative; it is therefore the truth set and
the formula is flat. The converse follows from flatness's principal
generator. For premise sets, select for each supported formula the first
supported resolution in the fixed ordering. This gives a supported
resolution of the set without an extra choice principle. Conversely,
support of a selected resolution supports its original formula, for
every premise.
\end{proof}
\end{annexitem}

\begin{annexitem}{annex:p8.inq.split}{thm:p08-inq-split}{Logical and contextual split and resolution}
\begin{proof}
A state supports a set \(\Gamma\) of flat formulas exactly when it is
contained in \(|\Gamma|_M\). Consequently the largest substate of \(s\)
supporting \(\Gamma\) is \(s\cap|\Gamma|_M\), and persistence proves
the specificity equivalence. The logical version takes the whole universe
in every model.

The contextual split property follows by evaluating the inquisitive
disjunction at that one largest state. For the logical version, suppose
\(\Gamma\not\models\varphi\) and \(\Gamma\not\models\psi\).
There are models with \(|\Gamma|\) failing the respective conclusions.
Take disjoint copies of these models and their union valuation. Locality
preserves each failure on its original component. The union of the two
\(\Gamma\)-truth sets supports \(\Gamma\) but cannot support either
conclusion: if it did, persistence would make the corresponding failing
component support it. It therefore fails the inquisitive disjunction.
This proves logical split by contraposition. Repeated finite splitting
and semantic normal form yield the logical and contextual resolution
properties for flat premises.

For arbitrary \(\Phi\), if \(\Phi\models\psi\), any
\(\Gamma\in R(\Phi)\) entails every member of \(\Phi\), so entails
\(\psi\). Since \(\Gamma\) is classical, the just-proved resolution
property supplies \(\beta\in R(\psi)\) with \(\Gamma\models\beta\).
Conversely, any state supporting \(\Phi\) supports a resolution
\(\Gamma\) of it. The assumed consequence from that resolution supplies
a supported resolution of \(\psi\), hence \(\psi\) itself. The same
argument restricted to substates of a fixed \(s\) proves the contextual
version. A dependence function is obtained by choosing the first successful
conclusion resolution, of which there are finitely many. Finally,
the split equivalence for a flat antecedent is the contextual split
property internalized by implication.
For a fixed single-antecedent map \(f\), support of \(\gamma_f\)
means support of every \(\alpha\to f(\alpha)\), equivalently
\(\alpha\models_s f(\alpha)\) for each antecedent resolution.
This is exactly the dependence-function condition. The semantic
resolution theorem for implication then makes support of the implication
equivalent to the existence of such a map.
\end{proof}
\end{annexitem}

\begin{annexitem}{annex:p8.inq.soundness}{thm:p08-inq-soundness}{Every primitive rule and the derived classical rules}
\begin{proof}
Induct on a proof, maintaining that any state supporting all undischarged
assumptions supports its conclusion. Conjunction introduction and either
elimination follow from the conjunction clause. Inquisitive disjunction
introduction follows from support of the chosen disjunct. Its elimination
is sound because the major premise supports one disjunct; the corresponding
branch induction hypothesis supplies the common conclusion.

Implication elimination uses the current state itself in the implication
clause. For introduction, any substate supporting the temporary antecedent
still supports all the other assumptions by persistence; the branch
induction hypothesis gives the consequent. This proves support of the
implication. Discharging zero or several occurrences changes none of the
argument: any retained occurrence is still among the undischarged assumptions.
Falsum elimination is sound because its premise forces the state to be
empty, and empty supports every formula. Double-negation elimination is
sound for classical \(\alpha\) because it is flat. Split is sound by
the proved flat-antecedent split equivalence. These exhaust the primitive
rules, so the induction proves soundness.

Negation introduction and elimination are the implication rules with
\(\bot\) as consequent. To introduce the defined classical disjunction
from \(\varphi\), assume \(\neg\varphi\land\neg\psi\), extract
\(\neg\varphi\), contradict \(\varphi\), and discharge that conjunction.
The result is \(\neg(\neg\varphi\land\neg\psi)\). The proof from
\(\psi\) is symmetric.

For classical disjunction elimination, suppose proofs from \(\Phi,\varphi\)
and from \(\Phi,\psi\) both conclude a classical \(\alpha\).
Given \(\Phi,\varphi\lor\psi\), temporarily assume \(\neg\alpha\).
Assuming \(\varphi\), use the first given proof and \(\neg\alpha\)
to derive \(\bot\), hence discharge \(\varphi\) to infer
\(\neg\varphi\). The second proof analogously gives \(\neg\psi\).
Their conjunction contradicts the defined disjunction. Discharge
\(\neg\alpha\), obtaining \(\neg\neg\alpha\), then apply the permitted
classical double-negation rule. No appeal to completeness is needed.

If arbitrary question conclusions were permitted in this last rule, the
two derivations \(p\vdash?p\) and \(\neg p\vdash?p\) would give
\(p\lor\neg p\vdash?p\). The two-world counterstate in the body
refutes this. It also supports \(\neg\neg?p\) but not \(?p\), refuting
unrestricted double-negation elimination.
\end{proof}
\end{annexitem}

\begin{annexitem}{annex:p8.inq.proof-resolution}{thm:p08-inq-proof-resolution}{The proof-to-resolution algorithm, rule by rule}
\begin{proof}
Label the finitely many assumption positions so that discharge is tracked
even when two formulas coincide. Unused assumptions may be retained as
weakening. Fix one resolution for each undischarged assumption. We define
the returned proof by induction on the last rule; the induction works
uniformly for the additional temporary assumptions in subproofs.

\paragraph{Assumption.}
Return its supplied resolution as an assumption. If the original proof
uses it several times, copy this one-line proof at each occurrence.

\paragraph{Conjunction introduction.}
The induction hypotheses return proofs of resolutions \(\beta\) and
\(\gamma\) of the two operands. Join them by conjunction introduction,
obtaining their conjunction, a resolution of the original conclusion.
\paragraph{Either conjunction elimination.}
The premise algorithm returns a conjunction of two operand resolutions.
Apply the matching projection to obtain the required one. These are
separate left and right cases with the same construction.

\paragraph{Implication introduction.}
Suppose the subproof concludes \(\xi\) under a temporary \(\chi\).
Enumerate \(R(\chi)=\{\beta_1,\ldots,\beta_m\}\).
For each \(i\), run the subproof algorithm with the fixed external
resolutions and the temporary resolution \(\beta_i\). It returns a
classical proof of some \(\gamma_i\in R(\xi)\). Discharge the
\(\beta_i\)-assumptions to obtain \(\beta_i\to\gamma_i\), and conjoin
the \(m\) results. This is the resolution \(\gamma_f\) of
\(\chi\to\xi\) for \(f(\beta_i)=\gamma_i\). If an assumption is unused,
the source permits its vacuous discharge.

\paragraph{Implication elimination.}
The antecedent subproof returns \(\beta\in R(\chi)\). The implication
subproof returns a conjunction with one conjunct
\(\beta_i\to\gamma_i\) for each \(\beta_i\in R(\chi)\).
Find the index with \(\beta=\beta_i\), project that conjunct, and apply
implication elimination to obtain \(\gamma_i\in R(\psi)\).
All the index sets here are the same finite resolution set; no independent
or mismatched enumeration is used.

\paragraph{Either inquisitive disjunction introduction.}
Return the proof provided by the operand algorithm unchanged. Its
conclusion belongs to the union of the two resolution sets.

\paragraph{Inquisitive disjunction elimination.}
Run the major-premise algorithm to obtain \(\beta\), a resolution of
one of its disjuncts. Choose that branch; if \(\beta\) is a resolution
of both, use a fixed tie convention. Run the branch algorithm with
\(\beta\) as the temporary resolution. It returns a classical proof
of some resolution \(\gamma\) of the desired conclusion. Substitute
copies of the major-premise output proof for every undischarged temporary
\(\beta\)-leaf in that branch proof. The resulting proof uses only
the external supplied resolutions and concludes \(\gamma\).
This proof substitution respects all original assumption labels.

\paragraph{Falsum elimination.}
The premise has the sole resolution \(\bot\), so its algorithm returns
a proof of \(\bot\). Choose the first member \(\beta\) of the finite
nonempty set \(R(\psi)\) and eliminate falsum to derive it. Return this
new classical proof, not the original unresolved proof.

\paragraph{Split.}
For classical \(\alpha\), its resolution set is \(\{\alpha\}\).
The resolutions of
\(\alpha\to(\varphi\mathbin{\vee\!\vee}\psi)\) are exactly
\(\alpha\to\beta\) for \(\beta\in R(\varphi)\cup R(\psi)\).
This is also precisely the resolution set of the split conclusion.
Return the premise algorithm's proof unchanged.

\paragraph{Classical double-negation elimination.}
The premise \(\neg\neg\alpha\) is classical, so is its own sole resolution.
The premise algorithm returns a classical proof of it. Apply classical
double-negation elimination to obtain \(\alpha\), the sole resolution
of the conclusion.

Every returned proof uses only classical formulas and classical
natural-deduction rules. The constructions call algorithms only for
proper subproofs, finitely many times, and all resolution sets are finite.
They therefore define a terminating procedure. Its output entailment is
sound by classical soundness, so the associated output-resolution map is
a logical dependence function. The empty premise list is included: its
unique resolution is the empty list.
\end{proof}
\end{annexitem}

\begin{annexitem}{annex:p8.inq.syntactic-resolutions}{thm:p08-inq-syntactic-resolutions}{Provable normal form and resolutions of arbitrary premise sets}
\begin{proof}
The rules other than split and classical double-negation elimination
include ordinary intuitionistic natural deduction with inquisitive
disjunction in its disjunction position. Its finite distributivity and
implication laws can therefore be derived here. To specify the needed
constructions, conjunction of two finite disjunctions distributes by
successively eliminating each disjunction and reintroducing the selected
conjunction as a disjunct; the reverse direction eliminates a conjunction
disjunct, projects its operands, and introduces each original disjunction.
An implication from a finite disjunction is interderivable with the
conjunction of the implications from its disjuncts: restrict the assumption
for one direction, and use disjunction elimination for the other.

Prove the resolution normal form by structural induction. Atoms and falsum
are immediate. For inquisitive disjunction, flatten the union of resolution
lists by its introduction/elimination rules. For conjunction, substitute
the operand normal forms and use the finite distribution just described.
For implication, substitute the normal forms of antecedent and consequent.
First distribute the antecedent disjunction into a conjunction of
implications. Every antecedent resolution is classical, so repeated split
turns each implication to a finite consequent disjunction into the
disjunction of its separate implications. The reverse of split is
derivable by assuming the common antecedent in each branch, applying
the chosen implication, introducing the consequent disjunction, and
discharging the antecedent. Finally distribute the finite conjunction of
these disjunctions. Each resulting disjunct chooses one consequent
resolution for each antecedent resolution, precisely the function
\(f\) and the formula \(\gamma_f\) in the definition of \(R\).
This proves interderivability for every constructor, without semantic
completeness. Every resolution therefore derives its original formula.

If every \(\alpha\in R(\varphi)\) together with \(\Phi\) derives
\(\psi\), repeated disjunction elimination and the proved normal form
give \(\Phi,\varphi\vdash\psi\). Contraposition gives the one-premise
failure assertion.

For a countably infinite premise set enumerate its members
\(\varphi_0,\varphi_1,\ldots\). Starting with failure of deduction,
repeatedly replace \(\varphi_n\) by a resolution \(\alpha_n\) while
retaining failure, using the preceding assertion and keeping all remaining
unresolved premises. The set of all selected \(\alpha_n\) cannot derive
\(\psi\): a finite derivation would use only finitely many of them,
contradicting the stage after they had all been selected. Finite premise
sets use the same argument for finitely many steps; the empty set needs
no choices. The countability follows from the fixed countable alphabet.

If \(\Phi\vdash\psi\), apply the proof-resolution algorithm to the
finitely many assumptions occurring in a proof and to the corresponding
members of any given resolution of \(\Phi\). It returns a derivable
resolution of \(\psi\). Conversely, if \(\Phi\not\vdash\psi\),
the selected classical resolution of \(\Phi\) just constructed cannot
derive any resolution of \(\psi\), since that would derive \(\psi\).
This proves the syntactic resolution equivalence. For classical premises
the sole resolution is that same premise set. A resolution of an inquisitive
disjunction belongs to one of the two operand resolution sets and derives
that operand, proving the final split assertion.
\end{proof}
\end{annexitem}

\begin{annexitem}{annex:p8.inq.support-lemma}{thm:p08-inq-support-lemma}{Classical theories and the full canonical support induction}
\begin{proof}
For consistency of classical formulas, the proof-resolution algorithm
turns any inquisitive derivation of \(\bot\) into a classical derivation
from exactly the same classical assumptions. Classical consistency and
consistency in this calculus consequently agree. Alternatively, the
usual completion argument can be performed directly: enumerate the
classical formulas, successively choose a formula or its negation while
retaining consistency, and take deductive closure. If both choices
were inconsistent, the corresponding classical negation reasoning and
classical double-negation elimination would make the preceding set
inconsistent. Finiteness of proofs preserves consistency at the union.
The closure is therefore a complete classical theory extending the input.

An intersection of theories is closed under classical deduction because
each deduction from the intersection is also a deduction from each
member theory. The empty intersection is all classical formulas and
is closed as well. In a classical theory, membership of a classical
formula is equivalent to its derivability.

We prove the support lemma by structural induction, simultaneously for
all states \(S\). For an atom, support means membership in every world,
equivalently membership in \(\bigcap S\), hence derivability. Falsum
is supported just when \(S=\varnothing\). In that case the intersection
contains \(\bot\). If \(S\) is nonempty and its intersection derived
\(\bot\), every world in \(S\) would derive it, contrary to consistency.
Conjunction follows from the two induction hypotheses and conjunction's
deduction rules.

For inquisitive disjunction, semantic support selects a supported disjunct;
induction and introduction derive the disjunction from the intersection.
Conversely, if the intersection derives it, syntactic split for classical
premises supplies a derivable disjunct. Induction makes that disjunct
supported, and hence the disjunction supported.

For implication, first suppose \(\bigcap S\vdash\varphi\to\psi\).
If \(T\subseteq S\) supports \(\varphi\), induction gives
\(\bigcap T\vdash\varphi\). Also \(\bigcap S\subseteq\bigcap T\),
so the latter theory derives the implication and then \(\psi\).
Induction yields \(T\models\psi\), proving implication support.

Conversely suppose \(\bigcap S\not\vdash\varphi\to\psi\).
Implication introduction implies \(\bigcap S,\varphi\not\vdash\psi\).
Traceable one-premise failure chooses \(\alpha\in R(\varphi)\) with
\(\bigcap S,\alpha\not\vdash\psi\). Let
\(T_\alpha=\{\Gamma\in S:\alpha\in\Gamma\}\).
The intersection of this substate contains \(\alpha\), including the
case of empty intersection, and therefore derives \(\varphi\).
By induction \(T_\alpha\models\varphi\).

If \(\bigcap T_\alpha\) derived \(\psi\), syntactic resolution for
classical premises would give some \(\beta\in R(\psi)\) belonging to
\(\bigcap T_\alpha\). Every \(\Gamma\in T_\alpha\) would then
contain the classical implication \(\alpha\to\beta\). Every other
\(\Gamma\in S\) contains \(\neg\alpha\), by completeness, and
therefore also contains \(\alpha\to\beta\). The elementary derivations
are respectively weakening followed by implication introduction, and
contradiction followed by falsum elimination and implication introduction.
Hence \(\alpha\to\beta\in\bigcap S\), making
\(\bigcap S,\alpha\vdash\beta\vdash\psi\), a contradiction.
Thus \(\bigcap T_\alpha\not\vdash\psi\), and induction gives
\(T_\alpha\not\models\psi\). This substate witnesses failure of the
implication at \(S\). In particular it is nonempty, since the empty
state supports every formula. All cases of the support induction are complete.
\end{proof}
\end{annexitem}

\begin{annexitem}{annex:p8.inq.completeness}{thm:p08-inq-completeness}{Canonical counterstates and compactness}
\begin{proof}
Soundness was proved independently. Suppose \(\Phi\not\vdash\psi\).
The syntactic resolution lemma supplies \(\Gamma\in R(\Phi)\) which
does not derive any member of \(R(\psi)=\{\alpha_1,\ldots,\alpha_n\}\).
For each \(i\), the classical set \(\Gamma\cup\{\neg\alpha_i\}\)
is consistent: otherwise negation introduction and classical
double-negation elimination would yield \(\Gamma\vdash\alpha_i\).
Extend it to a complete classical theory \(\Delta_i\), and take the
canonical state \(S=\{\Delta_1,\ldots,\Delta_n\}\).

Every world in \(S\) contains \(\Gamma\); the support lemma therefore
makes \(S\) support \(\Gamma\), hence \(\Phi\) by semantic resolution.
If \(S\models\psi\), semantic resolution gives some supported
\(\alpha_i\). The support lemma then gives
\(\bigcap S\vdash\alpha_i\), so \(\alpha_i\in\Delta_i\).
But \(\neg\alpha_i\in\Delta_i\) by construction, contradicting its
consistency. Thus \(S\) is a counterstate. It contains at most \(n\)
worlds, and \(n\geq1\). This proves completeness, including infinite
premise sets. Finally, every proof is finite and uses only finitely many
premises. Completeness and soundness therefore imply compactness of
the semantic entailment relation.
\end{proof}
\end{annexitem}

\begin{annexitem}{annex:p8.inq.partition-limit}{thm:p08-inq-partition-limit}{Failure of an internal implication in partition contents}
\begin{proof}
Suppose such a partition content \(Q\) existed. For each even world
\(e\) and odd world \(o\), take the unique-answer question whose
partition has the block \(\{e,o\}\) and has the four other worlds
as singleton blocks. This question together with parity resolves
outcome: intersecting any block with either parity class gives at
most one world. Thus it entails \(Q\), so \(e,o\) must be in one
block of the partition for \(Q\).

This is true for every even-odd pair. Fixing an odd world connects
all even worlds to its block, and fixing an even world connects all
odd worlds to that same block. Every world therefore belongs to
one block. In particular the presupposed domain cannot omit any
world, and \(Q\) is the trivial whole-universe partition content.
But parity itself entails this trivial content, while parity
together with parity still does not resolve outcome. That contradicts
the required equivalence for \(\lambda=\mathsf{parity}\).

For the inquisitive implication, a state supports parity-to-outcome
exactly when its even portion and odd portion each contain at most
one world. The maximal such states choose one even and one odd,
giving nine overlapping pairs. The implication/contextual consequence
lemma proves the requested internalization for this content.
\end{proof}
\end{annexitem}

\begin{annexitem}{annex:p8.inq.finite-calculations}{thm:p08-inq-finite-calculations}{Finite alternatives, entailments and resolution lists}
\begin{proof}
For the code table, a statement's unique alternative is its truth
set. The first four rows are therefore their stated singleton,
row, row-column union and material-conditional set. A polar question
has the two complementary truth-set alternatives, yielding the
fifth, sixth and eighth rows. The seventh row selects either of
the two stated row answers. Exact outcome selects one singleton;
first digit selects one row; a missing digit selects the complement
of its row-column union. The final conditional restricts to the
first row and then chooses one of its three singleton answers;
all worlds outside that row remain compatible with each alternative.
None of the alternatives in a given row properly contains another.
Equality, ordering and sum of two digits partition the nine worlds
exactly as the following three descriptions specify.

For the seven-column entailment table, the first row has the
singleton \(\{12\}\), which is included in an alternative of
every listed content. The row for the statement \(b\) has the
three-world set \(R_1\), contained in the alternatives for
\(b,c,f,h,j\), but not in the singleton for \(a\) or any
outcome singleton. The row for \(c\) has \(A_1\), contained
in alternatives only for \(c,h\) among the listed targets.
The question \(f\) has blocks \(R_1,W\setminus R_1\): its
second block mixes whether 1 occurs and mixes first digits, so
it entails none of \(h,i,j\), and its total presupposition entails
none of the non-tautological statements \(a,b,c\). The same
argument for \(h\), whose blocks are \(A_1,W\setminus A_1\),
shows it entails only itself among the targets. Each singleton
outcome resolves all four questions \(f,h,i,j\), while their
union is the whole space and entails none of \(a,b,c\).
Finally, each first-digit row resolves \(f\), but the second
and third rows mix the occurrence of 1 and have several outcomes;
thus \(j\) entails \(f,j\) and no other target. This checks
every entry, using the generator criterion rather than only truth
at the single worlds.

For the four-world table, \(p\land?q\) fixes \(p=1\) and
one value of \(q\), giving its two singletons. The inquisitive
disjunction of the two polar questions takes the union of their
four block alternatives. Negating the operand of a polar question
swaps its two blocks. The truth-value vector of \(p\lor q\)
and \(p\land q\) groups the worlds into
\(\{00\},\{01,10\},\{11\}\), proving the fourth row.
For \(\neg q\to?p\), the \(q=1\) worlds are unrestricted
and exactly one of the two \(q=0\) worlds may be added.
For \((p\lor q)\to(?p\land?q)\), world 00 is unrestricted
and at most one of the other three worlds may remain. The
equivalence of \(?p\) and \(?\neg p\) proves the seventh row.

For the last row, a state fails \((?p\to?q)\to?q\) precisely
when it contains a substate on which \(q\) varies but is determined
by \(p\). Such a substate contains a pair differing in both
coordinates, and conversely either diagonal pair is such a substate.
Hence support means avoiding both diagonal pairs. Its four maximal
states are exactly the two \(p\)-blocks and two \(q\)-blocks
listed in the table.

For the die implications, a classical antecedent \(E\) leaves
its complement \(O\) unrestricted, and restricts the live even
worlds to a consequent alternative. This yields the first three
listed results. The consequence from a range antecedent to the
flat statement even depends only on its tautological presupposition,
so is equivalent to even. For parity-to-range, each parity class
contains one world in each range block. The state may therefore
retain at most one even and one odd world, giving all nine maximal
pairs. For range-to-parity, each two-world range block contains
opposite parities, so at most one of each pair may remain; maximality
chooses one from each, giving eight triples. The same restriction
is obtained for range-to-outcome.

Finally, the four syntactic resolution computations apply respectively
the conjunction product, the implication function clause with one
antecedent resolution, the four functions between two two-element
resolution sets, and the product of two polar resolution sets.
These clauses do not delete a contradictory conjunction or identify
distinct but equivalent classical formulas. This proves the
resolution assertions exactly at their syntactic level.
\end{proof}
\end{annexitem}
